\chapter{Planificacion temporal}

\section{Descripcion de las tareas}
La planificación inicial estimaba 545 horas de dedicación distribuidas desde febrero hasta junio de 2026. El cómputo final asciende a 584 horas, una desviación de 39 horas (7,2\%) respecto a la previsión. El incremento no fue uniforme: se concentró en la implementación e integración, especialmente en la capa de voz STT/TTS, la compatibilidad con GPU AMD y la recuperación de servicios tras el cambio de red local.

\subsection{Gestion del proyecto}
Planificación, seguimiento, revisión de alcance, reuniones académicas y organización de la documentación.

\subsection{Analisis e investigacion}
Estudio de herramientas local-first, modelos LLM, motores de inferencia, seguridad perimetral, orquestadores, RAG y limitaciones del hardware.

\subsection{Diseno de la solucion}
Definición de topología de servicios, redes de contenedores, flujos de automatización, estructura de almacenamiento, acceso remoto y estrategia de copias de seguridad.

\subsection{Implementacion}
Despliegue de infraestructura, instalación de herramientas, integración de servicios, configuración de modelos locales, flujos de n8n, domótica, nube privada y monitorización.

\subsection{Pruebas y validacion}
Pruebas funcionales, latencia, estabilidad, uso de VRAM, recuperación ante fallos, seguridad de acceso remoto y validación de casos de uso reales.

\section{Estimacion de tiempos}
\begin{table}[H]
  \centering
  \caption{Estimacion de tiempos por fase}
  \label{tab:estimacion-tiempos}
  \begin{tabularx}{\textwidth}{lXr}
    \toprule
    Codigo & Tarea & Horas \\
    \midrule
    GP & Gestion y seguimiento del proyecto & 70 \\
    AN & Analisis e investigacion & 45 \\
    DS & Diseno de la solucion & 60 \\
    IM & Implementacion e integracion & 220 \\
    QA & Pruebas y validacion & 50 \\
    DO & Redaccion de memoria y defensa & 100 \\
    \bottomrule
  \end{tabularx}
\end{table}

\section{Hitos del proyecto}
\begin{itemize}
  \item \textbf{M1}: cierre de planificación, requisitos y alcance.
  \item \textbf{M2}: servidor físico operativo y contenedores base desplegados.
  \item \textbf{M3}: modelos LLM locales respondiendo sin dependencia externa.
  \item \textbf{M4}: integración mínima funcional entre voz, automatización y servicio privado.
  \item \textbf{M5}: memoria final, evidencias técnicas y defensa.
\end{itemize}

\section{Planificación inicial y ejecución real}
Para distinguir con claridad lo previsto de lo realizado se presentan dos diagramas. Ambos utilizan la misma escala semanal, desde febrero hasta junio de 2026. El primero representa el plan de partida; el segundo recoge la ejecución final y permite observar los solapamientos, retrasos y ampliaciones que aparecieron durante el desarrollo.

\begin{figure}[H]
  \centering
  \makebox[\textwidth][c]{%
  \resizebox{0.95\textwidth}{!}{%
  \begin{ganttchart}[
      hgrid,
      vgrid={*{3}{draw=none}, dotted},
      x unit=0.62cm,
      y unit title=0.6cm,
      y unit chart=0.7cm,
      title height=1,
      bar height=0.6,
      group height=0.3,
      title label font=\small,
      bar label font=\small,
      group label font=\small\bfseries,
      bar/.append style={fill=blue!40, rounded corners=2pt},
      bar incomplete/.append style={fill=gray!30},
      milestone/.append style={fill=orange, rounded corners=2pt}
    ]{1}{22}
    \gantttitle{Feb}{4}
    \gantttitle{Mar}{4}
    \gantttitle{Abr}{5}
    \gantttitle{May}{5}
    \gantttitle{Jun}{4} \\
    \ganttbar[bar/.append style={fill=gray!50}]{Gestión}{1}{22} \\
    \ganttbar[bar/.append style={fill=blue!30}]{Análisis}{1}{8} \\
    \ganttbar[bar/.append style={fill=teal!40}]{Diseño}{3}{10} \\
    \ganttbar[bar/.append style={fill=blue!55}]{Implementación}{6}{17} \\
    \ganttbar[bar/.append style={fill=violet!45}]{STT y TTS}{9}{13} \\
    \ganttbar[bar/.append style={fill=orange!50}]{Pruebas y validación}{14}{19} \\
    \ganttbar[bar/.append style={fill=green!40}]{Memoria y defensa}{10}{22} \\
    \ganttmilestone{M1}{4} \\
    \ganttmilestone{M2}{8} \\
    \ganttmilestone{M3}{12} \\
    \ganttmilestone{M4}{16} \\
    \ganttmilestone{M5}{22}
  \end{ganttchart}
  }
  }
  \caption{Diagrama de Gantt inicial del Proyecto Odín}
  \label{fig:gantt-inicial}
\end{figure}

\begin{figure}[H]
  \centering
  \makebox[\textwidth][c]{%
  \resizebox{0.95\textwidth}{!}{%
  \begin{ganttchart}[
      hgrid,
      vgrid={*{3}{draw=none}, dotted},
      x unit=0.62cm,
      y unit title=0.6cm,
      y unit chart=0.7cm,
      title height=1,
      bar height=0.6,
      title label font=\small,
      bar label font=\small,
      milestone label font=\small,
      bar/.append style={fill=blue!40, rounded corners=2pt},
      milestone/.append style={fill=orange, rounded corners=2pt}
    ]{1}{22}
    \gantttitle{Feb}{4}
    \gantttitle{Mar}{4}
    \gantttitle{Abr}{5}
    \gantttitle{May}{5}
    \gantttitle{Jun}{4} \\
    \ganttbar[bar/.append style={fill=gray!50}]{Gestión}{1}{22} \\
    \ganttbar[bar/.append style={fill=blue!30}]{Análisis}{1}{10} \\
    \ganttbar[bar/.append style={fill=teal!40}]{Diseño}{3}{11} \\
    \ganttbar[bar/.append style={fill=blue!55}]{Implementación}{6}{21} \\
    \ganttbar[bar/.append style={fill=violet!60}]{STT y TTS}{9}{19} \\
    \ganttbar[bar/.append style={fill=orange!50}]{Pruebas y validación}{12}{21} \\
    \ganttbar[bar/.append style={fill=green!40}]{Memoria y defensa}{10}{22} \\
    \ganttmilestone{M1}{4} \\
    \ganttmilestone{M2}{9} \\
    \ganttmilestone{M3}{13} \\
    \ganttmilestone{M4}{20} \\
    \ganttmilestone{M5}{22}
  \end{ganttchart}
  }
  }
  \caption{Diagrama de Gantt real del Proyecto Odín}
  \label{fig:gantt-real}
\end{figure}

La diferencia más visible aparece en la integración. El plan suponía que, una vez desplegados los contenedores, la voz quedaría resuelta mediante una combinación directa de Whisper y Piper. En la práctica fue necesario probar Whisper, faster-whisper, Wyoming Whisper y un ASR propio basado en ONNX para la entrada, además de Piper, Kokoro, Chatterbox y PocketTTS para la salida. También hubo que medir arranque en frío, latencia sostenida, precisión, naturalidad y consumo de memoria, y después adaptar los servicios a Home Assistant y Open WebUI.

Esta dificultad no es exclusiva del prototipo. Whisper ofrece una base robusta gracias a su entrenamiento a gran escala, pero su precisión varía con ruido, acento, micrófono y otras perturbaciones \citep{radford2023whisper,speechRobustBench2024}. En síntesis, la evaluación tampoco depende de una única métrica: deben equilibrarse inteligibilidad, naturalidad, prosodia, robustez y eficiencia \citep{tan2021neuraltts}. En un asistente local se añade además una restricción práctica: una voz de mayor calidad deja de ser útil si tarda demasiado en comenzar o compite con el LLM por la memoria disponible. Estas limitaciones explican que el bloque STT/TTS se prolongara desde abril hasta finales de mayo y que Piper se conservara como opción estable, aunque otras voces resultaran más naturales.

\section{Gestion del riesgo}
Los riesgos principales se concentran en la compatibilidad del stack AMD/ROCm, la saturación de VRAM al combinar LLM, RAG, voz y visión, la dependencia de un único nodo físico, la seguridad del acceso remoto y el crecimiento incontrolado del alcance. La mitigación combina fallback a CPU, carga bajo demanda de modelos, snapshots, copias de seguridad, segmentación de red y poda de funcionalidades secundarias si el calendario lo exige.

\section{Evolución respecto a la planificación inicial}
La planificación inicial subestimó el esfuerzo de integración entre servicios heterogéneos. En el seguimiento realizado el 15 de mayo de 2026 se habían acumulado 425 horas de trabajo, distribuidas según la Tabla~\ref{tab:horas-reales-intermedias}. Esta cifra constituye una fotografía histórica del proyecto, no el cómputo definitivo de cierre.

\begin{table}[H]
  \centering
  \caption{Horas acumuladas en el seguimiento intermedio}
  \label{tab:horas-reales-intermedias}
  \begin{tabular}{lr}
    \toprule
    Bloque de trabajo & Horas acumuladas \\
    \midrule
    Gestión del proyecto & 40 \\
    Análisis y requisitos & 60 \\
    Diseño de la solución & 55 \\
    Implementación e integración & 205 \\
    Pruebas y aseguramiento de calidad & 25 \\
    Documentación & 40 \\
    \midrule
    \textbf{Total} & \textbf{425} \\
    \bottomrule
  \end{tabular}
\end{table}

La desviación se concentró en compatibilidad de GPU, redes Docker, proxies, herramientas de Open WebUI, servicios de voz y recuperación tras el cambio de red local. Como respuesta, se redujo la prioridad de funciones accesorias, se separaron las pruebas por capas y se reservó el tramo final para estabilización, copias de seguridad, medición y documentación.

El cómputo total del proyecto al cierre de la memoria se sitúa en 584 horas. La Tabla~\ref{tab:horas-reales-finales} recoge el desglose final por bloque, comparado con la previsión inicial.

\begin{table}[H]
  \centering
  \caption{Horas finales reales frente a estimación inicial}
  \label{tab:horas-reales-finales}
  \begin{tabular}{lrrr}
    \toprule
    Bloque de trabajo & Estimado & Real & Desviación \\
    \midrule
    Gestión del proyecto         & 70  & 52  & $-18$ \\
    Análisis e investigación     & 45  & 58  & $+13$ \\
    Diseño de la solución        & 60  & 63  & $+3$  \\
    Implementación e integración & 220 & 275 & $+55$ \\
    Pruebas y validación         & 50  & 50  & $0$  \\
    Redacción de memoria         & 100 & 86  & $-14$ \\
    \midrule
    \textbf{Total} & \textbf{545} & \textbf{584} & \textbf{$+39$} \\
    \bottomrule
  \end{tabular}
\end{table}

Para hacer visible dónde se produjo el incremento, la Tabla~\ref{tab:desviacion-implementacion} detalla los principales trabajos de implementación. Las cifras son una reconstrucción razonada a partir del seguimiento del proyecto y de las incidencias documentadas.

\begin{table}[H]
  \centering
  \caption{Desglose de la desviación en implementación e integración}
  \label{tab:desviacion-implementacion}
  \small
  \begin{tabularx}{\textwidth}{Xrrr}
    \toprule
    Actividad & Previstas & Reales & Desviación \\
    \midrule
    Infraestructura, Docker y redes & 55 & 65 & $+10$ \\
    Núcleo conversacional, RAG y tools & 55 & 60 & $+5$ \\
    Domótica, cámaras y automatización & 45 & 50 & $+5$ \\
    Voz local: STT, TTS y Wyoming & 25 & 58 & $+33$ \\
    Monitorización, backup y seguridad & 40 & 42 & $+2$ \\
    \midrule
    \textbf{Total de implementación} & \textbf{220} & \textbf{275} & \textbf{$+55$} \\
    \bottomrule
  \end{tabularx}
\end{table}

La gestión consumió menos horas porque no existía coordinación de equipo. En cambio, el análisis aumentó por la investigación de AMD/ROCm, modelos locales y alternativas de voz. La mayor desviación se produjo en STT/TTS: cada motor requería instalar dependencias distintas, comprobar su compatibilidad con CPU o GPU, descargar voces y modelos, exponer una API, integrarla mediante Wyoming o HTTP y repetir las pruebas con audio real. Varias opciones funcionaban de forma aislada, pero no reunían simultáneamente baja latencia, naturalidad, pronunciación correcta en español y estabilidad. Por ello se cerró el alcance con Piper como TTS operativo, faster-whisper y Wyoming Whisper como soluciones maduras de STT, y los motores más ambiciosos como experimentación documentada.

\section{Revisión de riesgos durante la ejecución}
\begin{table}[H]
  \centering
  \caption{Riesgos revisados tras la implementación}
  \label{tab:riesgos-revisados}
  \begin{tabularx}{\textwidth}{XlX}
    \toprule
    Riesgo & Nivel & Respuesta adoptada o pendiente \\
    \midrule
    Incompatibilidades AMD/ROCm & Alto & Validación por servicio, versiones fijadas y alternativa CPU. \\
    Agotamiento de VRAM o memoria & Alto & Límites de recursos, menor concurrencia y reinicio controlado de Ollama. \\
    Fallo físico o pérdida de datos & Alto & NVMe montado por UUID, backup Restic cifrado, comprobación de integridad y restauración de muestra. \\
    Deuda operativa por configuraciones dispersas & Alto & Inventario, guía de réplica y documentación de dependencias. \\
    Dependencia de automatizaciones externas & Medio & Webhooks verificables y rutas locales cuando resulta viable. \\
    Exposición de secretos & Alto & Rotación, variables de entorno y ejemplos sanitizados. \\
    Crecimiento descontrolado del alcance & Medio & Priorización del núcleo y traslado de ampliaciones a trabajo futuro. \\
    \bottomrule
  \end{tabularx}
\end{table}
